Cuffless blood pressure (BP) monitoring is a key goal in wearable health, with photoplethysmography (PPG) as a promising sensing modality. However, PPG-based BP estimation has long suffered from the cross-subject generalisation challenge: due to individual differences in vascular elasticity and peripheral resistance, traditional machine learning models degrade severely under strict subject-level splits (XGBoost SBP MAE = 15.7 mmHg on 517 held-out subjects, failing all BHS grades).

This thesis proposes a few-shot personalised calibration method based on large language models (LLMs). The key idea is to exploit the in-context learning ability of LLMs: by providing the model with a small number of labelled calibration samples from the target subject, it adapts to that individual's physiological characteristics at inference time, without any fine-tuning. Two prompt designs are compared --- raw numeric features and semantic clinical descriptions that retain quantitative anchors. On the complete PulseDB dataset (5,159 subject records) with a strict 70/20/10 split, semantic description with 20 chronological calibration samples achieves an SBP MAE of 5.52 mmHg and DBP MAE of 2.93 mmHg, with DBP meeting the BHS Grade A clinical standard and SBP meeting Grade B --- outperforming XGBoost baselines (15.7/8.89 global; 7.79/4.18 with bias correction) and numeric LLM calibration (5.99/3.32).

The experiments reveal two key findings. First, the LLM provides its value without fine-tuning: both numeric and semantic few-shot calibration beat conventional bias correction on the large test set. Second, the advantage of semantic abstraction is most pronounced when calibration data is scarce --- at K=3 it outperforms numeric few-shot by 4.5 mmHg on SBP, a gap that closes as K grows. Drawing on SensorLM and TimeSRL, it is further shown that semantic abstraction must retain quantitative anchors --- purely qualitative summaries collapse (25.0/23.5).

\textbf{Keywords}: cuffless blood pressure; PPG; large language model; few-shot learning; personalised calibration; cross-subject generalisation
